Put \(N=M+1\). Since \(k=\lceil N(1-\alpha/2)\rceil\), the elementary equivalence \(\lceil x\rceil\le n\Longleftrightarrow x\le n\) for integral \(n\) gives \[k\le M\Longleftrightarrow N(1-\alpha/2)\le N-1\Longleftrightarrow \alpha\ge 2/(M+1).\tag{1}\] Fix the adverse arm \(a\). Under \(\texttt{def:disjoint-orbit-law}\), the arm-\(a\) outcome is \(K\) on target-supported arm-\(a\) orbits and zero on every other arm-\(a\) orbit, while every arm-\((1-a)\) outcome is zero. These prescriptions are consistent because every target-supported adverse-arm orbit has zero calibration mass. Complete the prototype with a common fixed covariate. The prototype is constant on every actual \(\mathcal G\)-orbit; alternatively, after uniform group randomization, left multiplication by any fixed \(g\in\mathcal G\) preserves the uniform group element. Hence its law satisfies \(\texttt{ass:finite-group-latent-invariance}\). All outcomes lie in \(\{0,K\}\subseteq[-K,K]\), so \(\texttt{def:invariant-laws}\) places the law in \(\mathcal P_{\mathcal G,K}\). Independent replication, the known randomization, and the stipulated auxiliary selector draws give the declared experiment. If \(q^{\mathrm{obs}}_{w,a}(O)>0\), then \(O\) is not target-supported, so its arm-\(a\) outcome and predictor are zero. Observability of the selector therefore makes every selected adverse-arm score zero. Scores in arm \(1-a\) are also zero. Because \(k\le M\), both order-statistic radii are zero, whence \[I_{0,w}=I_{1,w}=\{0\},\qquad \mathcal C_w=\{0\}.\tag{2}\] Target independence puts \(T^a\) in a target-supported orbit almost surely, so \(Y_{0T}(a)=K\) and \(Y_{0T}(1-a)=0\). Thus \(\tau_T=K\) for \(a=1\) and \(\tau_T=-K\) for \(a=0\). Since \(K>0\), (2) has conditional coverage zero. This proof uses the actual quotient \(\mathcal U/\mathcal G\); in a pair example the cell \(O_{b,a}\) may be substituted only under the local condition in the corrected \(\texttt{def:pair-action}\).